%% Replace the author metadata marked TODO before circulation.
\documentclass[11pt]{article}

\usepackage{amsthm,amsmath,amsfonts,amssymb,mathtools}
\usepackage[numbers,sort&compress]{natbib}
\usepackage[hidelinks]{hyperref}
\usepackage{booktabs}
\usepackage[nameinlink,noabbrev]{cleveref}

\numberwithin{equation}{section}

\theoremstyle{plain}
\newtheorem{theorem}{Theorem}[section]
\newtheorem{proposition}[theorem]{Proposition}
\newtheorem{lemma}[theorem]{Lemma}
\newtheorem{corollary}[theorem]{Corollary}

\theoremstyle{definition}
\newtheorem{definition}[theorem]{Definition}
\newtheorem{remark}[theorem]{Remark}

\newcommand{\E}{\mathbb E}
\newcommand{\Pp}{\mathbb P}
\newcommand{\1}{\mathbf 1}
\newcommand{\sech}{\operatorname{sech}}
\newcommand{\RS}{\mathrm{RS}}
\newcommand{\AT}{\mathrm{AT}}
\newcommand{\cP}{\mathcal P}
\newcommand{\cK}{\mathcal K}
\newcommand{\cR}{\mathcal R}
\newcommand{\dd}{\,\mathrm d}
\newcommand{\norm}[1]{\left\lVert #1\right\rVert}
\newcommand{\ip}[2]{\left\langle #1,#2\right\rangle}
\newcommand{\Tr}{\operatorname{Tr}}
\newcommand{\diag}{\operatorname{diag}}
\newcommand{\sgn}{\operatorname{sgn}}
\newcommand{\Var}{\operatorname{Var}}
\newcommand{\Cov}{\operatorname{Cov}}
\newcommand{\e}{\mathrm e}

\title{A quantitative replica-symmetric bound in the strict AT region}
%% TODO: replace the following author and e-mail data.
\author{First Author\\
Department of Mathematics, Keio University\\
\texttt{author@example.edu}}
\date{}

\begin{document}
\maketitle

\begin{abstract}
For the Sherrington--Kirkpatrick model with deterministic external field
$h>0$, we prove a finite-volume replica-symmetric estimate, uniform on
compact subsets of the strict de Almeida--Thouless region and uniform
along the replica-symmetric smart path:
\[
 \E\big\langle(R_{12}-q)^2\big\rangle=O(N^{-1}).
\]
The argument does not invoke the identity between the limiting free
energy and the Parisi variational formula.  The principal ingredient is
a two-replica smart-path bound whose overlap--multiplier mixed
derivative is the replicon gap.  It yields a quadratic constrained
free-energy estimate.  A finite-volume Gronwall argument first gives
vanishing overlap fluctuations, Gaussian concentration upgrades this
to fixed-deviation exponential tails, and a three-dimensional cavity
system then gives the optimal second-moment bound.  We also obtain an
$O(N^{-1})$ replica-symmetric free-energy correction and identify the
finite-volume replicon susceptibility.
\end{abstract}

\medskip
\noindent\textit{Mathematics Subject Classification.}
Primary 60K35, 82B44; secondary 60G15, 60F10.

\smallskip
\noindent\textit{Keywords.}
Sherrington--Kirkpatrick model; de Almeida--Thouless condition;
Guerra--Talagrand interpolation; cavity method; overlap concentration.

\section{Introduction, model, and main results}\label{sec:intro}

\subsection{Introduction}

The Sherrington--Kirkpatrick model is the basic mean-field model of a
disordered Ising system.  In a nonzero external field, its
replica-symmetric order parameter is determined by a scalar fixed-point
equation, while the de Almeida--Thouless condition describes the
stability of that solution against two-replica perturbations.  The
strict inequality in the AT condition is therefore expected to imply
that the overlap has fluctuations of order $N^{-1/2}$ and that the
replica-symmetric free energy is accurate up to order $N^{-1}$.

The aim of this paper is to prove these finite-volume conclusions
uniformly on compact subsets of the strict AT region.  The proof is
organized so that the stability mechanism is visible at every step.
The constrained two-replica pressure has a flat replica-symmetric
saddle at overlap $q$.  Introducing both an overlap-path parameter and
a Lagrange multiplier isolates the mixed derivative at that saddle;
the derivative is exactly the negative replicon gap
$-(1-s\alpha)$.  This yields a quadratic constrained-pressure loss.
A finite-volume Gronwall argument then gives preliminary overlap
concentration, which is upgraded by Gaussian concentration to
exponentially small fixed-deviation tails.  Finally, a
three-dimensional cavity system converts the fixed-scale information
into the optimal second-moment estimate.

The smart-path and cavity computations are classical tools in the SK
model; see \cite{Talagrand,Panchenko}.  The general two-dimensional
Guerra--Talagrand variational bound was developed in \cite{ChenGT}, and
related finite-size fluctuation systems appear in \cite{DeyWu}.  The
strict-AT diffusion comparison used below is closely related to
\cite{Lopatto}.  For completeness, we prove in the present
replica-symmetric finite-step setting every interpolation identity,
PDE derivative, and cavity equation needed for the argument.  We use
without proof only standard probability tools: Gaussian integration
by parts and concentration, It\^o's formula, conditional expectation,
H\"older and Cauchy--Schwarz inequalities, and Gronwall's inequality.
In particular, the proof does not use the identity between the
limiting free energy and the Parisi variational formula.

\subsection{Model and the strict AT region}

Let $\Sigma_N=\{-1,1\}^N$.  For $\beta>0$ and $h>0$, let
\[
 H_N(\sigma)
 =
 \frac{\beta}{\sqrt N}\sum_{1\le i<j\le N}g_{ij}\sigma_i\sigma_j
 +h\sum_{i=1}^N\sigma_i ,
\]
where the $g_{ij}$ are independent standard normal random variables.
For replicas $\sigma^1,\sigma^2,\ldots$ sampled independently from the
Gibbs measure for fixed disorder, put
\[
 R_{ab}=\frac1N\sum_{i=1}^N\sigma_i^a\sigma_i^b.
\]

Let $q=q(\beta,h)$ be the unique solution of
\begin{equation}
 q=\E\tanh^2\bigl(h+\beta\sqrt q\,Z\bigr),
 \qquad Z\sim N(0,1).
 \label{eq:q}
\end{equation}
Existence and uniqueness for $h>0$ are standard; see, for example,
\cite[Proposition~1.3.8]{Talagrand}.  Continuity follows directly from
uniqueness: if $(\beta_n,h_n)\to(\beta,h)$, every convergent
subsequence of $q(\beta_n,h_n)\in[0,1]$ has, by dominated convergence
in \eqref{eq:q}, a limit solving the fixed-point equation at
$(\beta,h)$.  Uniqueness identifies that limit with $q(\beta,h)$, so
the entire sequence converges.
Set
\begin{equation}
 r=\E\tanh^4\bigl(h+\beta\sqrt q\,Z\bigr),\qquad
 a=1-2q+r=\E\sech^4\bigl(h+\beta\sqrt q\,Z\bigr),
 \label{eq:ra}
\end{equation}
and let
\begin{equation}
 \alpha=\beta^2a.
 \label{eq:alpha}
\end{equation}
We work on a fixed compact set
\begin{equation}
 \cK\Subset\{(\beta,h):\beta>0,\ h>0,\ \alpha(\beta,h)<1\}.
 \label{eq:K}
\end{equation}
Thus, for some $\delta_{\cK}>0$,
\begin{equation}
 1-\alpha\ge\delta_{\cK}
 \qquad\text{on }\cK.
 \label{eq:delta}
\end{equation}
The maps $q,r,a,\alpha$ are therefore continuous on $\cK$.  In
particular, compactness and $h>0$ imply
\begin{equation}
 0<q_{\cK}:=\inf_{(\beta,h)\in\cK}q
 \le \sup_{(\beta,h)\in\cK}q<1.
 \label{eq:qcompact}
\end{equation}
Indeed, $q=0$ would contradict \eqref{eq:q} and $h>0$, while
$q<1$ because $\tanh^2(h+\beta\sqrt q\,Z)<1$ almost surely.

\subsection{The replica-symmetric smart path}

For $0\le s\le1$, define the smart-path Hamiltonian
\begin{equation}
 H_{N,s}(\sigma)
 =
 \frac{\beta\sqrt s}{\sqrt N}
 \sum_{1\le i<j\le N}g_{ij}\sigma_i\sigma_j
 +\sum_{i=1}^N
 \bigl(h+\beta\sqrt{(1-s)q}\,z_i\bigr)\sigma_i,
 \label{eq:path}
\end{equation}
where $(z_i)$ is an independent standard Gaussian family.  We write
$\langle\cdot\rangle_s$ for the Gibbs expectation and
\[
 \nu_s(f)=\E\langle f\rangle_s,\qquad Q_{ab}=R_{ab}-q.
\]
We also write
\begin{equation}
 \phi_N(s)=\frac1N\E\log
 \sum_{\sigma\in\Sigma_N}e^{H_{N,s}(\sigma)}.
 \label{eq:pathfreeenergy}
\end{equation}
The effective interaction strength is $\beta\sqrt s$.  At the
replica-symmetric overlap $q$, the interaction and random-field
variances add to
\begin{equation}
 s\beta^2q+(1-s)\beta^2q=\beta^2q.
 \label{eq:constantvariance}
\end{equation}
Hence the one-site effective field has the same law
$h+\beta\sqrt q\,Z$ for every $s$, while the AT parameter is
\begin{equation}
 \alpha_s=s\alpha,\qquad 1-\alpha_s\ge\delta_{\cK}.
 \label{eq:alphas}
\end{equation}

\subsection{Main result}

Define
\begin{equation}
 A_s=\nu_s(Q_{12}^2),\qquad
 B_s=\nu_s(Q_{12}Q_{13}),\qquad
 C_s=\nu_s(Q_{12}Q_{34}).
 \label{eq:ABC}
\end{equation}

\begin{theorem}[Quantitative strict-AT bound]
\label{thm:main}
There is a constant $C_{\cK}<\infty$ such that, for every $N\ge1$,
\[
 \sup_{\substack{(\beta,h)\in\cK\\0\le s\le1}}
 N\,\nu_s(Q_{12}^2)\le C_{\cK}.
\]
Consequently, if
\[
 \phi_N(\beta,h)=\frac1N\E\log\sum_{\sigma\in\Sigma_N}e^{H_N(\sigma)}
\]
and
\[
 \Phi_{\RS}(\beta,h)
 =
 \log2+\E\log\cosh(h+\beta\sqrt q\,Z)
 +\frac{\beta^2}{4}(1-q)^2,
\]
then
\[
 0\le\Phi_{\RS}(\beta,h)-\phi_N(\beta,h)
 \le\frac{C_{\cK}}N
\]
uniformly on $\cK$.

Moreover, uniformly for $(\beta,h,s)\in\cK\times[0,1]$,
\[
 N\bigl(A_s-2B_s+C_s\bigr)
 =
 \frac{a}{1-s\alpha}+o_{\cK}(1).
\]
Here and below, $o_{\cK}(1)\to0$ as $N\to\infty$, uniformly over the
displayed compact parameter set.
\end{theorem}

\subsection{Proof strategy}

The proof is divided into the following steps.  First, the scalar
finite-step PDE is solved explicitly, and the strict AT sign is proved
uniformly along the smart path.  Second, a specialized two-replica
Guerra--Talagrand interpolation is derived from a finite cascade and
Gaussian interpolation.  The strict AT sign then yields a uniform
quadratic bound for the constrained pressure.  Third, a quadratically
coupled pressure and Gronwall's inequality produce an $o(1)$ overlap
bound, and Gaussian concentration gives fixed-deviation exponential
tails.  Fourth, a cavity expansion gives a three-dimensional linear
system for $A_s,B_s,C_s$ with a cubic remainder.  Uniform invertibility
of the cavity matrix and the fixed-deviation estimate allow that
remainder to be absorbed.  The free-energy estimate follows by
integrating the exact smart-path identity, and the replicon
susceptibility follows from the replicon row of the same cavity system.

\section{Proof of the main results}\label{sec:proof}

\subsection{Gaussian calculus used throughout}\label{subsec:gaussian}

We first record the two elementary Gaussian identities repeatedly used
below.  Let $G=(G_1,\ldots,G_m)$ be centered Gaussian with covariance
matrix $C$, and let $F$ be continuously differentiable with bounded
first derivatives.  Gaussian integration by parts gives
\begin{equation}
 \E[G_iF(G)]=\sum_{j=1}^m C_{ij}\E[\partial_jF(G)].
 \label{eq:GIBP}
\end{equation}
If $G^{(0)}$ and $G^{(1)}$ are independent centered Gaussian vectors
with covariance matrices $C_0$ and $C_1$, respectively, and
$G_t=\sqrt t\,G^{(1)}+\sqrt{1-t}\,G^{(0)}$, then applying
\eqref{eq:GIBP} twice yields, initially for $0<t<1$,
\begin{equation}
 \frac{\dd}{\dd t}\E F(G_t)
 =\frac12\sum_{i,j}(C_1-C_0)_{ij}
   \E[\partial_{ij}F(G_t)].
 \label{eq:gaussianinterpolation}
\end{equation}
The right-hand side has a continuous extension to $t=0,1$ whenever
the displayed expectations do.  Both formulas extend to the log
partition functions below by a standard truncation argument, since all
spin observables are bounded.

We shall also use the following immediate consequence of Gaussian
concentration.  If $F_v$, $v\in I$, are functions of the same standard
Gaussian vector, each with Lipschitz constant at most $L$, and $I$ is
finite, then
\begin{equation}
 \E\max_{v\in I}\{F_v-\E F_v\}
 \le L\sqrt{2\log|I|}.
 \label{eq:gaussianmaxgeneral}
\end{equation}
Indeed, the exponential moment bound
$\E\exp\{t(F_v-\E F_v)\}\le\exp(t^2L^2/2)$, followed by the union
bound and optimization in $t$, proves \eqref{eq:gaussianmaxgeneral}.

\subsection{The scalar replica-symmetric PDE and the strict AT sign}\label{subsec:ATsign}

For fixed $(\beta,h,s)\in\cK\times[0,1]$, put
\begin{equation}
 b=\beta\sqrt s,\qquad d=\beta^2(1-s)q,
 \qquad x_q(u)=\1_{[q,1]}(u),\qquad
 \xi_s(u)=du+\frac{b^2}{2}u^2.
 \label{eq:bdxq}
\end{equation}
Let $\Psi:[0,1]\times\mathbb R\to\mathbb R$ solve
\begin{equation}
 \partial_u\Psi+\frac{b^2}{2}
 \bigl(\partial_{xx}\Psi+x_q(u)(\partial_x\Psi)^2\bigr)=0,
 \qquad \Psi(1,x)=\log\cosh x.
 \label{eq:parisipde}
\end{equation}
This finite-step equation has the explicit semigroup representation
\begin{align}
 \Psi(u,x)
 &=\log\cosh x+\frac{b^2}{2}(1-u), &&q\le u\le1,
 \label{eq:PDEupper}\\
 \Psi(u,x)
 &=\E\Psi\bigl(q,x+b\sqrt{q-u}\,Z\bigr), &&0\le u\le q.
 \label{eq:PDElower}
\end{align}
Indeed, on $[q,1]$ the Cole--Hopf transform $\exp\Psi$ solves the
backward heat equation, and on $[0,q]$ the equation itself is the
backward heat equation.  These formulas prove existence, uniqueness,
and all differentiability statements used below.

Let $Z_0$ and a Brownian motion $W$ be independent and define the local
field process by
\begin{equation}
 \dd X_u=b^2x_q(u)\partial_x\Psi(u,X_u)\,\dd u+b\,\dd W_u,
 \qquad X_0=h+\sqrt d\,Z_0.
 \label{eq:localfield}
\end{equation}
The replica-symmetric trial value along the path is
\begin{equation}
 P_s^*=\log2+\E\Psi(0,h+\sqrt d\,Z_0)
 -\frac{b^2}{2}\int_q^1u\,\dd u.
 \label{eq:scalartrial}
\end{equation}
Using \eqref{eq:PDEupper}--\eqref{eq:PDElower} and
$d+b^2q=\beta^2q$, this reduces to
\begin{equation}
 P_s^*=\log2+\E\log\cosh(h+\beta\sqrt q\,Z)
 +\frac{s\beta^2}{4}(1-q)^2.
 \label{eq:RSpathvalue}
\end{equation}


\begin{proposition}[Uniform strict AT sign]
\label{prop:pathRS}
Let $\Psi=\Psi_{\delta_q}$ and $X=X^{\delta_q}$, and set
\[
 g_s(u)=\E\bigl[\partial_x\Psi(u,X_u)^2\bigr].
\]
Uniformly on $\cK\times[0,1]$,
\begin{equation}
 g_s(u)-u
 \begin{cases}
 >0,&0\le u<q,\\
 =0,&u=q,\\
 <0,&q<u\le1.
 \end{cases}
 \label{eq:ATsign}
\end{equation}
Moreover, there are $c_{\cK},\varepsilon_{\cK}>0$ such that
\begin{equation}
 |g_s(u)-u|\ge c_{\cK}|u-q|
 \qquad\text{whenever }|u-q|\le\varepsilon_{\cK}.
 \label{eq:linearATsign}
\end{equation}
The sign is uniform away from $q$ on the compact family.
\end{proposition}

\begin{proof}
For $u\ge q$,
$\partial_x\Psi(u,x)=\tanh x$, while for $u\le q$,
\[
 \partial_x\Psi(u,X_u)=\E[\tanh(X_q)\mid\mathcal F_u],
 \qquad
 \partial_{xx}\Psi(u,X_u)=\E[\sech^2(X_q)\mid\mathcal F_u].
\]
The initial Gaussian belongs to $\mathcal F_u$, and
\begin{equation}
 X_q\stackrel d=h+\beta\sqrt q\,Z
 \label{eq:Xq}
\end{equation}
by \eqref{eq:constantvariance}.  Ito's formula and conditional Jensen
therefore give, for $u\le q$,
\[
 g_s'(u)
 =b^2\E[\partial_{xx}\Psi(u,X_u)^2]
 \le b^2\E\sech^4(X_q)=s\alpha.
\]
Since $g_s(q)=q$,
\begin{equation}
 g_s(u)-u\ge(1-s\alpha)(q-u)
 \qquad(0\le u<q).
 \label{eq:leftstrict}
\end{equation}

For $u\ge q$, put $m_u=\tanh X_u$.  The drift cancels:
\begin{equation}
 \dd m_u=b(1-m_u^2)\,\dd W_u.
 \label{eq:mdiffusion}
\end{equation}
Thus, if $f(u)=\E m_u^2-u$,
\begin{equation}
 f(q)=0,\qquad f'(u)=b^2\E(1-m_u^2)^2-1,
 \qquad f'(q)=s\alpha-1.
 \label{eq:fprime}
\end{equation}
If $b^2(1-q)\le1$, then
$f'(q)=s\alpha-1\le-\delta_{\cK}<0$, so $f<0$ immediately to the
right of $q$.  If $f$ ever returned to zero, let $u_0>q$ be its first
zero after this initial negative interval.  Then differentiability
would give $f'(u_0)\ge0$.  On the other hand,
\[
 \E(1-m_{u_0}^2)^2\le1-\E m_{u_0}^2=1-u_0,
\]
so $f'(u_0)\le b^2(1-u_0)-1
<b^2(1-q)-1\le0$, a contradiction.  Hence $f(u)<0$ for every
$u>q$.

Suppose $b^2(1-q)>1$ and put $\sigma^2=\beta^2q$.  We first
show that $h<\sigma^2$.  The function
$c\mapsto\E\sech^2(c+\sigma Z)$ is nonincreasing on $[0,\infty)$.
Indeed, if $f(x)=\sech^2x$ and $\varphi_\sigma$ is the centered Gaussian
density, then for $c\ge0$,
\[
 \frac{\dd}{\dd c}(f*\varphi_\sigma)(c)
 =\int_0^\infty f'(x)
   \{\varphi_\sigma(x-c)-\varphi_\sigma(x+c)\}\,\dd x\le0,
\]
because $f'(x)\le0$ for $x\ge0$ and the expression in braces is
nonnegative.  Hence, if $h\ge\sigma^2$, then
\[
 1-q=\E\sech^2(h+\sigma Z)
 \le\E\sech^2(\sigma^2+\sigma Z).
\]
Let $\varepsilon$ be a symmetric $\{-1,1\}$-valued random variable
and observe $Y=\sigma\varepsilon+Z$.  Its conditional mean is
$\E[\varepsilon\mid Y]=\tanh(\sigma Y)$.  Therefore
\[
 \E\sech^2(\sigma^2+\sigma Z)
 =\E\bigl(\varepsilon-\E[\varepsilon\mid Y]\bigr)^2.
\]
Conditional expectation minimizes mean-square error, while the best
linear estimator of $\varepsilon$ from $Y$ has error
$(1+\sigma^2)^{-1}$.  Consequently
\[
 1-q\le(1+\sigma^2)^{-1},
\]
Multiplying by $1+\beta^2q$ and using $q>0$ gives
$\beta^2q(1-q)\le q$, hence $\beta^2(1-q)\le1$, a contradiction.  Thus $h<\sigma^2$.

Use the time variable $\lambda=b^2(u-q)$ and set
$v_\lambda=1-m_{q+\lambda/b^2}^2$.  Conditional on $X_q=x$,
Girsanov's formula gives
\[
 \E[\psi(X_{q+\lambda/b^2})\mid X_q=x]
 =
 \frac{\E[\psi(x+\sqrt\lambda G)\cosh(x+\sqrt\lambda G)]}
      {\E[\cosh(x+\sqrt\lambda G)]}.
\]
Writing $S=\sech^2X_q$, take $\psi=\sech^4$ in the preceding
transition formula and use the symmetry of $G$.  Conditional on
$X_q=x$, this gives
\[
 \E[v_\lambda^2\mid X_q=x]
 =\frac{\e^{-\lambda/2}\sech x}{2}\E_G
 \bigl[\sech^3(x+\sqrt\lambda G)
       +\sech^3(x-\sqrt\lambda G)\bigr].
\]
For $y\in\mathbb R$ and $S=\sech^2x$, elementary algebra yields
\[
 \frac{\sech x}{2}\{\sech^3(x+y)+\sech^3(x-y)\}
 =S^2\cosh y\,
 \frac{1+(4-3S)\sinh^2y}{(1+S\sinh^2y)^3}.
\]
Averaging first in $G$ and then in $X_q$ therefore gives
\begin{equation}
 \E v_\lambda^2=\E[S^2F_\lambda(S)],
 \label{eq:Flambda-representation}
\end{equation}
where
\begin{equation}
 F_\lambda(y)=\e^{-\lambda/2}\E\left[
 \cosh(\sqrt\lambda G)
 \frac{1+(4-3y)\sinh^2(\sqrt\lambda G)}
      {(1+y\sinh^2(\sqrt\lambda G))^3}\right].
 \label{eq:Flambda}
\end{equation}
We verify the two monotonicity facts needed for
\eqref{eq:Flambda-representation}.  Put $t=\sinh^2(\sqrt\lambda G)$
and
\[
 H_t(y)=\frac{1+(4-3y)t}{(1+yt)^3}.
\]
Then
\[
 H_t'(y)=-\frac{6t\{1+(2-y)t\}}{(1+yt)^4}\le0,
\]
so $F_\lambda$ is decreasing.  Moreover, with
$y=1-u^2$,
\[
 \int_0^1H_t(y)\frac{\dd y}{2\sqrt{1-y}}
 =\int_0^1\frac{1+(1+3u^2)t}{(1+(1-u^2)t)^3}\,\dd u
 =\left[\frac{u}{(1+t-tu^2)^2}\right]_{0}^{1}=1.
\]
Since $\e^{-\lambda/2}\E\cosh(\sqrt\lambda G)=1$, it follows that
\begin{equation}
 \int_0^1F_\lambda(y)\frac{\dd y}{2\sqrt{1-y}}=1.
 \label{eq:Fmeanone}
\end{equation}

It remains to compare the law of $S$ with the reference probability
measure $\mu_0(\dd y)=\dd y/(2\sqrt{1-y})$.  Write
$x(y)=\operatorname{arcosh}(y^{-1/2})$.  Since
$X_q\sim N(h,\sigma^2)$, the density of the finite measure
$\E[S^2\1_{\{S\in\dd y\}}]$ with respect to $\mu_0$ is, up to a
positive constant,
\begin{equation}
 \rho(y)=y\exp\left\{-\frac{x(y)^2+h^2}{2\sigma^2}\right\}
 \cosh\left(\frac{h x(y)}{\sigma^2}\right).
 \label{eq:rhodensity}
\end{equation}
As a function of $x\ge0$, its logarithmic derivative equals
\[
 -2\tanh x-\frac{x}{\sigma^2}
 +\frac{h}{\sigma^2}\tanh\left(\frac{hx}{\sigma^2}\right).
\]
When $h\le\sigma^2$, the last term is at most $\tanh x$; hence this
derivative is strictly negative.  Since $y=\sech^2x$ decreases in
$x$, $\rho$ is increasing in $y$.  Therefore $F_\lambda$ and $\rho$
have opposite monotonicities.  Applying
\[
 \int fg\,\dd\mu_0\le\left(\int f\,\dd\mu_0\right)
 \left(\int g\,\dd\mu_0\right)
\]
for a decreasing $f$ and an increasing $g$---an identity obtained by
integrating $(f(x)-f(y))(g(x)-g(y))\le0$ over $\mu_0^{\otimes2}$---and
using \eqref{eq:Fmeanone}, we obtain
\begin{equation}
 \E v_\lambda^2\le\E S^2=a.
 \label{eq:strictATdiffusion}
\end{equation}
It follows from \eqref{eq:fprime} that
$f'(u)\le s\alpha-1<0$, proving the right-hand sign.

Finally, \eqref{eq:leftstrict} and
$f'(q)=-(1-s\alpha)\le-\delta_{\cK}$ give
\eqref{eq:linearATsign} by uniform continuity.  Compactness gives
uniform strictness away from $q$.  This completes the proof.
\end{proof}

\subsection{Two-replica Guerra--Talagrand coercivity}\label{subsec:GT}

For an attainable overlap $v\in\mathcal R_N
:=\{-1,-1+2/N,\ldots,1\}$, define
\begin{equation}
 Z^{(2)}_{N,s}(v)
 =
 \sum_{\substack{\sigma^1,\sigma^2\in\Sigma_N\\R_{12}=v}}
 e^{H_{N,s}(\sigma^1)+H_{N,s}(\sigma^2)}
 \label{eq:constrainedZ}
\end{equation}
and retain the notation $P_s^*$ from \eqref{eq:RSpathvalue}.  For
$\lambda\in\mathbb R$, define
\begin{equation}
 f_\lambda(x_1,x_2)
 =\log\left(\frac14\sum_{\varepsilon_1,\varepsilon_2=\pm1}
 \exp\{\varepsilon_1x_1+\varepsilon_2x_2
       +\lambda\varepsilon_1\varepsilon_2\}\right).
 \label{eq:GTterminal}
\end{equation}

\begin{lemma}[Finite-step two-replica smart-path bound]
\label{lem:specialGT}
Assume that the centered one-replica Gaussian Hamiltonian has covariance
$N\xi_s(R(\sigma,\tau))$.  Let
$Q:[0,1]\to\mathbb S_+^2$ be a continuous piecewise $C^1$ path, let
$\gamma:[0,1]\to[0,1]$ be a nondecreasing step function, and
assume
\begin{equation}
 Q_0=0,\qquad Q_1=
 \begin{pmatrix}1&v\\v&1\end{pmatrix},
 \qquad
 A(u)=\frac{\dd}{\dd u}\xi_s'(Q_u)\succeq0
 \label{eq:GTlemmaassumptions}
\end{equation}
where scalar functions are applied entrywise.  Let $U$ solve
\[
 \partial_uU+\frac12\bigl(\ip{A(u)}{D^2U}
 +\gamma(u)\ip{A(u)\nabla U}{\nabla U}\bigr)=0
\]
with terminal condition $f_\lambda$ from \eqref{eq:GTterminal}.  If the
self-overlap matrix of the constrained pair is $Q_1$, then
\begin{align}
 \frac1N\E\log Z^{(2)}_{N,s}(v)
 \le{}&2\log2+\E U(0,Y_0,Y_0)-\lambda v\nonumber\\
 &-\frac12\int_0^1\gamma(u)
 \sum_{a,b=1}^2A_{ab}(u)Q_{u,ab}\,\dd u.
 \label{eq:specialGTlemma}
\end{align}
Here the off-diagonal entry of $Q_1$ is $v$ and
$Y_0=h+\sqrt d\,Z_0$.
\end{lemma}

\begin{proof}
We give the finite-cascade construction in detail.  First suppose
that all values of $\gamma$ lie in $(0,1)$.  Choose a partition
\begin{equation}
 0=u_0<u_1<\cdots<u_K=1
 \label{eq:GTpartition}
\end{equation}
containing every discontinuity of $\gamma$ and every break point of
$Q$, and write $m_k$ for the value of $\gamma$ on
$[u_{k-1},u_k)$.  Equal consecutive values can be merged.

For $0<m<1$, let $(u_j)_{j\ge1}$ be the atoms, ordered decreasingly,
of a Poisson point process on $(0,\infty)$ with intensity
$m x^{-1-m}\,\dd x$.  If $(X_j)$ are i.i.d. and independent of the
normalized weights $w_j=u_j/\sum_\ell u_\ell$, the scaling property of
the Poisson process gives
\begin{equation}
 \E\log\sum_jw_j\e^{X_j}
 =\frac1m\log\E\e^{mX_1}.
 \label{eq:PDidentity}
\end{equation}
Indeed, if $S=\sum_j u_j$, the marking theorem shows that
$(u_je^{X_j})$ is a Poisson process with the original intensity
multiplied by $\E e^{mX_1}$.  Hence
$\sum_j u_je^{X_j}$ has the same law as
$(\E e^{mX_1})^{1/m}S$.  Taking the difference of the expected
logarithms of these two sums proves \eqref{eq:PDidentity}.
More explicitly, at every vertex
$\alpha^{k-1}=(\alpha_1,\ldots,\alpha_{k-1})$ take an independent
copy $(u_{\alpha^{k-1}j})_{j\ge1}$ of the Poisson process with
parameter $m_k$, with its atoms ordered decreasingly.  For a leaf
$\alpha=(\alpha_1,\ldots,\alpha_K)$ form the raw product
\[
 v_\alpha=\prod_{k=1}^K u_{\alpha^{k-1}\alpha_k}
\]
and normalize the leaf masses by
$w_\alpha=v_\alpha/\sum_{\alpha'}v_{\alpha'}$.  This is the finite
Ruelle cascade.  The finiteness needed for normalization follows
successively from $0<m_1<\cdots<m_K<1$; the inequalities are strict
because equal consecutive values were merged.  If a terminal variable
$X_K$ is built from independent increments at the $K$ levels, set
\begin{equation}
 X_{k-1}=\frac1{m_k}\log\E_k e^{m_kX_k},
 \qquad 1\le k\le K,
 \label{eq:cascaderecursion}
\end{equation}
where $\E_k$ averages only the level-$k$ increment.  Conditioning
from the leaves back to the root and applying
\eqref{eq:PDidentity} at each vertex gives
\begin{equation}
 \E\log\sum_\alpha w_\alpha e^{X_K(\alpha)}=\E X_0.
 \label{eq:cascadeidentity}
\end{equation}

We now construct the Gaussian fields.  Put, entrywise,
\begin{equation}
 C_k=\xi_s'(Q_{u_k}),\qquad
 \Theta(x)=x\xi_s'(x)-\xi_s(x),\qquad
 D_k=\sum_{a,b=1}^2\Theta(Q_{u_k,ab}).
 \label{eq:CDlevels}
\end{equation}
The assumptions imply
\begin{equation}
 C_k-C_{k-1}
 =\int_{u_{k-1}}^{u_k}A(u)\,\dd u\succeq0.
 \label{eq:Cincrements}
\end{equation}
Furthermore, $D_0=0$, and differentiation on each smooth piece gives
\begin{equation}
 D_k-D_{k-1}
 =\int_{u_{k-1}}^{u_k}
 \sum_{a,b=1}^2A_{ab}(u)Q_{u,ab}\,\dd u\ge0.
 \label{eq:Dincrements}
\end{equation}
Indeed, the integrand is $\Tr(A(u)Q_u)\ge0$.

For every site $i$, take a base Gaussian vector $z_{i,0}$ with
covariance $C_0$, and on every branch independent Gaussian increments
$z_{i,k}$ with covariance $C_k-C_{k-1}$.  Set
\[
 z_i(\alpha)=z_{i,0}+\sum_{k=1}^K
 z_{i,k}(\alpha_1,\ldots,\alpha_k).
\]
Then
\begin{equation}
 \E z_i^a(\alpha)z_i^b(\alpha')
 =\xi_s'\bigl(Q_{u_{\alpha\wedge\alpha'},ab}\bigr).
 \label{eq:zcovariance}
\end{equation}
Since $Q_0=0$ and $\xi_s'(0)=d$,
$z_{i,0}=(\sqrt d\,Z_i,\sqrt d\,Z_i)$.  Similarly, independent
scalar branch increments with variances $D_k-D_{k-1}$ define
$y(\alpha)$ with
\begin{equation}
 \E y(\alpha)y(\alpha')=D_{\alpha\wedge\alpha'}.
 \label{eq:ycovariance}
\end{equation}
Thus \eqref{eq:Cincrements}--\eqref{eq:Dincrements} explicitly prove
existence of both hierarchical fields, including singular increments.

For a constrained pair
$\boldsymbol\sigma=(\sigma^1,\sigma^2)$, define
\begin{align}
 \mathcal H_t(\boldsymbol\sigma,\alpha)
 ={}&\sqrt t\,[H_{N,s}^{\rm cen}(\sigma^1)
                     +H_{N,s}^{\rm cen}(\sigma^2)]\nonumber\\
 &+\sqrt{1-t}\sum_{i=1}^N\sum_{a=1}^2
 z_i^a(\alpha)\sigma_i^a+\sqrt{tN}\,y(\alpha)\nonumber\\
 &+h\sum_{i=1}^N(\sigma_i^1+\sigma_i^2)
 +\lambda\sum_{i=1}^N\sigma_i^1\sigma_i^2-\lambda Nv.
 \label{eq:GTinterpolatingHamiltonian}
\end{align}
Here $H_{N,s}^{\rm cen}$ is the full centered Gaussian part of
\eqref{eq:path}.  The last two $\lambda$-terms cancel on
$R_{12}=v$.

Let $\E_{\mathrm{cas}}$ denote expectation over the cascade weights
and all hierarchical Gaussian increments, and set
\begin{equation}
 \varphi_N(t)=\frac1N\E\E_{\mathrm{cas}}\log
 \sum_{\substack{\boldsymbol\sigma\\R_{12}=v}}
 \sum_\alpha w_\alpha e^{\mathcal H_t(\boldsymbol\sigma,\alpha)}.
 \label{eq:GTphi}
\end{equation}
At $t=0$ we remove the constraint:
\begin{align}
 &\sum_{\substack{\boldsymbol\sigma\\R_{12}=v}}
 e^{\sum_{i,a}z_i^a(\alpha)\sigma_i^a
 +h\sum_{i,a}\sigma_i^a
 +\lambda\sum_i\sigma_i^1\sigma_i^2-\lambda Nv}\nonumber\\
 &\quad\le e^{-\lambda Nv}
 \prod_{i=1}^N\sum_{\varepsilon_1,\varepsilon_2=\pm1}
 e^{\varepsilon_1(h+z_i^1(\alpha))
   +\varepsilon_2(h+z_i^2(\alpha))
   +\lambda\varepsilon_1\varepsilon_2}.
 \label{eq:GTremoveconstraint}
\end{align}
The logarithm of each one-site factor is
$2\log2+f_\lambda(h+z_i^1,h+z_i^2)$.  Since branch increments are
independent between sites, \eqref{eq:cascaderecursion} factorizes at
every level.  On a segment where $\gamma=m$, the Cole--Hopf transform
$e^{mU}$ solves the backward heat equation with covariance speed
$A$; its Gaussian semigroup is exactly
\eqref{eq:cascaderecursion}.  Iterating over the partition and
averaging the base vector gives
\begin{equation}
 \varphi_N(0)
 \le2\log2+\E U(0,Y_0,Y_0)-\lambda v.
 \label{eq:GTendpointzero}
\end{equation}

At $t=1$, the spin Hamiltonian is independent of $\alpha$.  For the
level-$k$ increment $G_k$ of $y$,
\[
 \frac1{m_k}\log\E_k
 e^{m_k(x+\sqrt N\,G_k)}
 =x+\frac{Nm_k}{2}(D_k-D_{k-1}).
\]
Consequently, \eqref{eq:Dincrements} gives
\begin{align}
 \varphi_N(1)
 ={}&\frac1N\E\log Z^{(2)}_{N,s}(v)\nonumber\\
 &+\frac12\int_0^1\gamma(u)
 \sum_{a,b=1}^2A_{ab}(u)Q_{u,ab}\,\dd u.
 \label{eq:GTendpointone}
\end{align}

For the interpolation derivative, take two samples
$(\boldsymbol\sigma,\alpha)$ and
$(\boldsymbol\tau,\alpha')$ and put
\[
 R^{ab}=\frac1N\sum_i\sigma_i^a\tau_i^b,\qquad
 Q^{ab}=Q_{u_{\alpha\wedge\alpha'},ab}.
\]
The covariance of the centered Gaussian part of
\eqref{eq:GTinterpolatingHamiltonian}, divided by $N$, is
\[
 t\sum_{a,b}\xi_s(R^{ab})
 +(1-t)\sum_{a,b}\xi_s'(Q^{ab})R^{ab}
 +t\sum_{a,b}\Theta(Q^{ab}).
\]
Its derivative is
\begin{align}
 \sum_{a,b}\{&
 \xi_s(R^{ab})-\xi_s'(Q^{ab})R^{ab}
 +Q^{ab}\xi_s'(Q^{ab})-\xi_s(Q^{ab})\}\nonumber\\
 &=\sum_{a,b}\Delta_{\xi_s}(R^{ab},Q^{ab}),
 \label{eq:GTcovariancederivative}
\end{align}
where
\begin{equation}
 \Delta_{\xi_s}(x,y)
 =\xi_s(x)-\xi_s(y)-\xi_s'(y)(x-y)
 =\frac{b^2}{2}(x-y)^2\ge0.
 \label{eq:GTDelta}
\end{equation}
We now justify the passage from finite to infinite cascades.  For
$M\ge1$, let
\[
 \mathcal A_M=\{\alpha:1\le\alpha_k\le M, 1\le k\le K\},
 \qquad V_M=\sum_{\alpha\in\mathcal A_M}v_\alpha,
 \qquad w_{\alpha,M}=\frac{v_\alpha}{V_M},
\]
and let $\varphi_{N,M}$ be defined as in \eqref{eq:GTphi}, with
$\alpha\in\mathcal A_M$ and weights $w_{\alpha,M}$.  The sets
$\mathcal A_M$ increase to the full leaf set and
\begin{equation}
 V_M\uparrow V:=\sum_\alpha v_\alpha\in(0,\infty)
 \quad\text{almost surely}.
 \label{eq:cascadetruncationmass}
\end{equation}
For completeness, this follows recursively from the marking theorem.
If $(u_j)$ has parameter $m$ and the independent nonnegative marks
$S_j$ satisfy $\E S_1^m<\infty$, then $\sum_j u_jS_j<\infty$ almost
surely.  At adjacent cascade levels the required moment is finite
because $m_{k-1}<m_k$; induction from the last level proves
\eqref{eq:cascadetruncationmass}.

Conditional on the truncated weights, the sum over
$(\boldsymbol\sigma,\alpha)$ is a finite Gaussian log partition
function.  Formula \eqref{eq:gaussianinterpolation} therefore applies
for $0<t<1$.  The resulting covariance formula extends to $t=0,1$
by continuity.  For the self term both matrices equal $Q_1$, so
\eqref{eq:GTDelta} vanishes, and for every $M$,
\begin{equation}
 \varphi_{N,M}'(t)
 =-\frac12\E\Big\langle
 \sum_{a,b=1}^2\Delta_{\xi_s}(R^{ab},Q^{ab})
 \Big\rangle_{t,M}\le0,
 \qquad 0<t<1.
 \label{eq:GTderivative}
\end{equation}
In particular,
$\varphi_{N,M}(1)\le\varphi_{N,M}(0)$.

It remains to pass only this endpoint inequality to the limit.  For
$t\in\{0,1\}$, write $Z_M(t)$ for the partition sum inside the
logarithm defining $\varphi_{N,M}(t)$ and $Z(t)$ for its full-cascade
counterpart.  Since $w_\alpha=v_\alpha/V$,
\[
 Z_M(t)=\frac{V}{V_M}
 \sum_{\alpha\in\mathcal A_M}w_\alpha
 \sum_{\boldsymbol\sigma:R_{12}=v}
 e^{\mathcal H_t(\boldsymbol\sigma,\alpha)}.
\]
The second factor increases to $Z(t)$, while $V/V_M\to1$.  The limit
is finite almost surely: conditional on the cascade weights, Gaussian
exponential moments and the finiteness of the spin set give
\[
 \E_{\rm G} Z(t)\le 4^N\exp(C_{N,s,\lambda})<\infty.
\]
Thus $\log Z_M(t)\to\log Z(t)$ almost surely.

These logarithms are uniformly integrable.  Indeed, conditional on
the weights, each is a Gaussian log-sum-exp whose squared Lipschitz
constant is at most $C_{\cK,\lambda}N$, independently of $M$.
Jensen's inequality gives a uniform upper bound $C_{\cK,\lambda}N$
for its Gaussian mean.  For the lower bound, retain one admissible
spin pair and use
$\log\sum_\alpha w_{\alpha,M}e^{X_\alpha}
 \ge\sum_\alpha w_{\alpha,M}X_\alpha$; after Gaussian averaging this
is bounded below by $-C_{\cK,\lambda}N$.  Gaussian concentration now
gives a uniform second-moment bound for $\log Z_M(t)$.  Consequently,
\begin{equation}
 \varphi_{N,M}(t)\longrightarrow\varphi_N(t),
 \qquad t=0,1.
 \label{eq:cascadeendpointlimit}
\end{equation}
Passing $M\to\infty$ in the endpoint inequality yields
$\varphi_N(1)\le\varphi_N(0)$.  Together with
\eqref{eq:GTendpointzero} and \eqref{eq:GTendpointone}, this proves the
lemma when $\gamma$ takes values in $(0,1)$.

For a general step function with values in $[0,1]$, apply the result
to
\[
 \gamma_\eta=(1-2\eta)\gamma+\eta,\qquad
 0<\eta<\frac12.
\]
The recursion operators
\[
 T_mF=m^{-1}\log\E e^{mF},\qquad T_0F=\E F,
\]
are continuous for $m\in[0,1]$.  Since the terminal function has at
most linear spatial growth, Gaussian exponential integrability gives
convergence of every finite recursion and of the correction integral
as $\eta\downarrow0$.  This proves the cases $m=0$ and $m=1$ without
postulating Poisson--Dirichlet weights at either endpoint.
\end{proof}

For $(\beta,h,s)\in\cK\times[0,1]$, retain
$b^2=s\beta^2$, $d=\beta^2(1-s)q$, and
$\xi_s(u)=du+b^2u^2/2$.  Given $v\in[-1,1]$, write $r_v=|v|$ and put
$\iota_v=1$ for $v\ge0$ and $\iota_v=-1$ for $v<0$.  Define
\begin{align}
 Q^v_u
 &=
 \begin{cases}
  \begin{pmatrix}u&\iota_vu\\ \iota_vu&u\end{pmatrix},
       &0\le u\le r_v,\\[2mm]
  \begin{pmatrix}u&v\\v&u\end{pmatrix},
       &r_v\le u\le1,
 \end{cases}
 \label{eq:GTpath}\\
 \gamma_v(u)
 &=
 \begin{cases}
  \frac12x_q(u),&0\le u<r_v,\\
  x_q(u),&r_v\le u\le1.
 \end{cases}
 \label{eq:halfparameter}
\end{align}
Both $Q^v$ and
\begin{equation}
 A^v_s(u):=\frac{\dd}{\dd u}\xi_s'(Q^v_u)
 =
 \begin{cases}
 b^2\begin{pmatrix}1&\iota_v\\\iota_v&1\end{pmatrix},&u<r_v,\\[2mm]
 b^2I_2,&u>r_v
 \end{cases}
 \label{eq:GTspeed}
\end{equation}
are positive semidefinite.  Scalar functions are applied entrywise to
matrices.  Let $U^{\lambda,v}_s$ be the finite-semigroup solution of
\begin{equation}
 \partial_uU+\frac12\left(
   \langle A^v_s,D^2U\rangle
   +\gamma_v(u)\langle A^v_s\nabla U,\nabla U\rangle
 \right)=0
 \label{eq:2DParisiPDE}
\end{equation}
with terminal condition $f_\lambda$ from \eqref{eq:GTterminal}.  With
$Y_0=h+\sqrt d\,Z_0$, set
\begin{align}
 \mathcal L_s(v)
 &=\frac12\int_0^1\gamma_v(u)
   \sum_{a,b=1}^2 A^v_{s,ab}(u)Q^v_{u,ab}\,\dd u,
 \label{eq:GTcorrectiondef}\\
 \mathfrak P_s(\lambda,v)
 &=2\log2+\E U^{\lambda,v}_s(0,Y_0,Y_0)
   -\lambda v-\mathcal L_s(v).
 \label{eq:GTfunctional}
\end{align}
Below $r_v$, all four summands in \eqref{eq:GTcorrectiondef} equal
$b^2u$; above $r_v$, only the two diagonal summands remain.  Hence
\begin{equation}
 \mathcal L_s(v)
 =b^2\int_0^1x_q(u)u\,\dd u
 =\frac{b^2}{2}(1-q^2),
 \label{eq:GTcorrection}
\end{equation}
independently of $v$.

\begin{lemma}[Continuity of the finite-semigroup functional]
\label{lem:GTcontinuity}
On every compact set on which $\lambda$ is bounded, the maps
\[
 (\beta,h,s,\lambda,v)\longmapsto
 \mathfrak P_s(\lambda,v)-2P_s^*,
 \qquad
 (\beta,h,s,\lambda,v)\longmapsto
 \partial_\lambda\mathfrak P_s(\lambda,v)
\]
are jointly continuous for
$(\beta,h,s,v)\in\cK\times[0,1]\times[-1,1]$.
\end{lemma}

\begin{proof}
Split $[0,1]$ at $q$ and $|v|$.  On each resulting interval the
solution is obtained by one of the three Gaussian recursion operators
\[
 \mathcal T_mF(x)=m^{-1}\log\E e^{mF(x+G)},
 \qquad m\in\{1/2,1\},
 \qquad
 \mathcal T_0F(x)=\E F(x+G),
\]
where the covariance of $G$ is the integral of $A_s^v$ over that
interval.  These covariances depend continuously on the parameters.
At $v=0$, at $v=\pm q$, or when another interval collapses, the
corresponding covariance tends to zero and its operator tends to the
identity.  The same holds at $s=0$, since every $b^2$ covariance
vanishes.  The terminal function has uniformly linear growth on
bounded $\lambda$-sets, so Gaussian exponential integrability and
dominated convergence prove continuity of the finite composition.
For the multiplier derivative, differentiate the terminal function
and propagate it through the same tilted Gaussian expectations; its
absolute value is at most one.  Dominated convergence applies again.
Finally, $q$, $d$, $Y_0$, \eqref{eq:GTcorrection}, and $P_s^*$ depend
continuously on the parameters.  This proves both assertions.
\end{proof}

\begin{proposition}[Uniform quadratic GT coercivity]
\label{prop:GTcoercivity}
There is a constant $c_{\cK}>0$ such that
\begin{equation}
 \frac1N\E\log Z^{(2)}_{N,s}(v)
 \le 2P_s^*-c_{\cK}(v-q)^2
 \label{eq:GTcoercivity}
\end{equation}
for all $(\beta,h,s)\in\cK\times[0,1]$, all $N$, and all
$v\in\mathcal R_N$.
\end{proposition}

\begin{proof}
We give the normalization in full.  This is important here: with a
different factor in either the two-dimensional cumulative parameter
or the Lagrange term, the mixed derivative below would not be the AT
replicon.

Recall from \eqref{eq:bdxq} that
$b^2=s\beta^2$, $d=\beta^2(1-s)q$, and
$\xi_s(u)=du+b^2u^2/2$.
The covariance of the centered Gaussian part of \eqref{eq:path} is
\begin{equation}
 \E H_{N,s}^{\rm cen}(\sigma)H_{N,s}^{\rm cen}(\tau)
 =N\xi_s(R(\sigma,\tau))-\frac{b^2}{2}.
 \label{eq:GTcovariance}
\end{equation}
If $G\sim N(0,b^2/2)$ is independent and is added to every
configuration, the covariance becomes exactly $N\xi_s(R)$.  The
one-replica log partition function is shifted by $G$, and the
two-replica log partition function by $2G$; their disorder
expectations are unchanged.  Define
\[
 \widetilde H_{N,s}^{\rm cen}(\sigma)
 =H_{N,s}^{\rm cen}(\sigma)+G.
\]
For the remainder of this application we denote
$\widetilde H_{N,s}^{\rm cen}$ again by $H_{N,s}^{\rm cen}$.  Thus its
covariance is exactly $N\xi_s(R)$, and the normalization used in
Lemma~\ref{lem:specialGT} introduces no finite-$N$ error.

Lemma~\ref{lem:specialGT}, applied after the harmless covariance normalization above, gives
\begin{equation}
 \frac1N\E\log Z^{(2)}_{N,s}(v)
 \le \mathfrak P_s(\lambda,v),
 \qquad \lambda\in\mathbb R.
 \label{eq:GTfunctionalbound}
\end{equation}
The statement is finite-volume and includes the rank-one pieces of the
path; no limiting argument in $N$ is involved.

We next identify the zero-source derivatives.  Suppose first that
$v\ge0$.  On $[v,1]$, the PDE at $\lambda=0$ factorizes as
\[
 U^{0,v}_s(u,x_1,x_2)=\Psi(u,x_1)+\Psi(u,x_2).
\]
On $[0,v]$, set $V(u,x)=U^{0,v}_s(u,x,x)$.  From
\eqref{eq:2DParisiPDE},
\[
 \partial_uV+\frac{b^2}{2}
 \left(V_{xx}+\frac{x_q(u)}2V_x^2\right)=0,
 \qquad V(v,x)=2\Psi(v,x).
\]
Hence $V=2\Psi$.  Combining this identity with
\eqref{eq:GTcorrection} and the scalar trial value
\eqref{eq:scalartrial} gives the exact flatness identity
\begin{equation}
 \mathfrak P_s(0,v)=2P_s^*,
 \qquad 0\le v\le1.
 \label{eq:GTflat}
\end{equation}
We now derive the multiplier derivative explicitly.  Set
\[
 W^{v}_s(u,x_1,x_2)
 =\left.\partial_\lambda
 U^{\lambda,v}_s(u,x_1,x_2)\right|_{\lambda=0}.
\]
Direct differentiation of the terminal function gives
\begin{equation}
 f_0(x_1,x_2)=\log\cosh x_1+\log\cosh x_2,\qquad
 \left.\partial_\lambda f_\lambda(x_1,x_2)\right|_{\lambda=0}
 =\tanh x_1\tanh x_2.
 \label{eq:GTterminalderivative}
\end{equation}
On $[v,1]$, $A_s^v=b^2I_2$ and
$U_s^{0,v}=\Psi(x_1)+\Psi(x_2)$.  Differentiating
\eqref{eq:2DParisiPDE} gives
\begin{align}
 \partial_uW_s^v+\frac{b^2}{2}\bigl\{
 \partial_{11}W_s^v+\partial_{22}W_s^v
 +2x_q(u)[\Psi_x(u,x_1)\partial_1W_s^v
          +\Psi_x(u,x_2)\partial_2W_s^v]\bigr\}=0.
 \label{eq:GTlinearizedupper}
\end{align}
Differentiating the scalar PDE in $x$ verifies that the solution with
terminal data \eqref{eq:GTterminalderivative} is
\begin{equation}
 W_s^v(u,x_1,x_2)
 =\Psi_x(u,x_1)\Psi_x(u,x_2),\qquad v\le u\le1.
 \label{eq:GTlinearizedproduct}
\end{equation}
On $[0,v]$, put $\overline W(u,x)=W_s^v(u,x,x)$.  Since
$A_s^v=b^2\begin{psmallmatrix}1&1\\1&1\end{psmallmatrix}$ and
$U_s^{0,v}(u,x,x)=2\Psi(u,x)$, the restricted linearized equation is
\begin{equation}
 \partial_u\overline W+\frac{b^2}{2}\overline W_{xx}
 +b^2x_q(u)\Psi_x(u,x)\overline W_x=0,\qquad
 \overline W(v,x)=\Psi_x(v,x)^2.
 \label{eq:GTlinearizedlower}
\end{equation}
The backward Kolmogorov formula for \eqref{eq:localfield} yields
\[
 \E W_s^v(0,Y_0,Y_0)=\E\Psi_x(v,X_v)^2=g_s(v).
\]
The correction $\mathcal L_s(v)$ is independent of $\lambda$, while
$\partial_\lambda(-\lambda v)=-v$.  We have therefore proved
\begin{equation}
 \partial_\lambda\mathfrak P_s(0,v)=g_s(v)-v.
 \label{eq:GTfirstderivative}
\end{equation}
In particular,
\begin{equation}
 \partial_v\partial_\lambda\mathfrak P_s(0,q)
 =b^2\E[\partial_{xx}\Psi(q,X_q)^2]-1
 =s\alpha-1.
 \label{eq:GTmixedderivative}
\end{equation}
This is the replicon mixed derivative with all normalizing factors
displayed.

The differentiated PDE also gives constants
$\lambda_{\cK},M_{\cK}>0$ such that
\begin{equation}
 \sup_{\substack{(\beta,h,s)\in\cK\times[0,1]\\
                  -1\le v\le1,\ |\lambda|\le\lambda_{\cK}}}
 |\partial_{\lambda\lambda}\mathfrak P_s(\lambda,v)|
 \le M_{\cK}.
 \label{eq:GTsecondderivative}
\end{equation}
This bound follows directly from the finite-step semigroup formula.
For a Gaussian vector $G$ and $m>0$, write
\[
 \mathcal T_mF(x)=m^{-1}\log\E\exp\{mF(x+G)\},
 \qquad \mathcal T_0F(x)=\E F(x+G).
\]
Then
\begin{align*}
 \partial_\lambda\mathcal T_mF
 &=\E_m[\partial_\lambda F],\\
 \partial_{\lambda\lambda}\mathcal T_mF
 &=\E_m[\partial_{\lambda\lambda}F]
   +m\Var_m(\partial_\lambda F),
\end{align*}
where $\E_m$ denotes the exponentially tilted expectation.  At the
terminal time, $|\partial_\lambda f_\lambda|\le1$ and
$0\le\partial_{\lambda\lambda}f_\lambda\le1$.  For every signed path
$-1\le v\le1$, the only positive semigroup parameters are $1/2$ and
$1$; splitting at $q$ and $|v|$ therefore creates at most two
nonzero levels.  The zero level is an ordinary Gaussian expectation.
Iterating the displayed identities gives
$|\partial_\lambda U|\le1$ and, for example,
$0\le\partial_{\lambda\lambda}U\le3$.  Averaging the base Gaussian and
differentiating $-\lambda v$ do not change the second derivative.
This proves \eqref{eq:GTsecondderivative} on the full signed range,
uniformly in all parameters.

By \eqref{eq:linearATsign} and
\eqref{eq:GTfirstderivative}, if $|v-q|$ is sufficiently small, then
\[
 |\partial_\lambda\mathfrak P_s(0,v)|
 \ge c_{\cK}|v-q|,
 \qquad
 \operatorname{sgn}\partial_\lambda\mathfrak P_s(0,v)
 =-\operatorname{sgn}(v-q).
\]
After decreasing that neighborhood so that
$|\partial_\lambda\mathfrak P_s(0,v)|/M_{\cK}\le\lambda_{\cK}$,
Taylor's theorem with
\[
 \lambda(v)=-\frac{
 \partial_\lambda\mathfrak P_s(0,v)}{M_{\cK}}
\]
yields
\begin{equation}
 \mathfrak P_s(\lambda(v),v)
 \le2P_s^*-\frac{
 |\partial_\lambda\mathfrak P_s(0,v)|^2}{2M_{\cK}}
 \le2P_s^*-c_{\cK}(v-q)^2.
 \label{eq:localGTgap}
\end{equation}

We next compute the signed derivative by the same semigroups.  Suppose
$-q\le v\le q$, put $r=|v|$ and
$\iota=\operatorname{sgn}(v)$, with $\iota=1$ when $v=0$.  The
rank-one interval $[0,r]$ lies below the first mass $q$.  At
$\lambda=0$, the $m=1$ recursion on $[q,1]$ and the diagonal heat
recursion on $[r,q]$ give
\[
 U_s^{0,v}(r,x_1,x_2)=\Psi(r,x_1)+\Psi(r,x_2).
\]
The remaining $m=0$ recursion is linear.  Starting it from
$(Y_0,Y_0)$ gives the pair
\[
 X_1=Y_0+b\sqrt r\,Z,\qquad
 X_2=Y_0+\iota b\sqrt r\,Z
\]
at time $r$.  Each marginal is the scalar heat evolution, so
\[
 \E U_s^{0,v}(0,Y_0,Y_0)=2\E\Psi(0,Y_0).
\]
Together with \eqref{eq:GTcorrection}, this proves the flatness
identity in \eqref{eq:signedflat}.

For the derivative, \eqref{eq:GTterminalderivative} and
\eqref{eq:GTlinearizedproduct} give
$W_s^v(r,x_1,x_2)=\Psi_x(r,x_1)\Psi_x(r,x_2)$.  The explicit scalar
heat semigroup gives
\[
 \Psi_x(r,x)=\E_G\tanh(x+b\sqrt{q-r}\,G).
\]
Introduce independent $G_1,G_2$, also independent of $(Y_0,Z)$, and
put
\[
 Y_1=Y_0+b\sqrt r\,Z+b\sqrt{q-r}\,G_1,\qquad
 Y_2=Y_0+\iota b\sqrt r\,Z+b\sqrt{q-r}\,G_2.
\]
Then both means are $h$, both variances are $d+b^2q=\beta^2q$, and
their covariance is $d+\iota b^2r=d+b^2v$.  It follows that
\[
 \E W_s^v(0,Y_0,Y_0)=\E\tanh(Y_1)\tanh(Y_2)=f_s(v).
\]
After differentiating $-\lambda v$, we have proved, rather than
assumed,
\begin{equation}
 \mathfrak P_s(0,v)=2P_s^*,\qquad
 \partial_\lambda\mathfrak P_s(0,v)=f_s(v)-v,
 \label{eq:signedflat}
\end{equation}
where
\begin{equation}
 f_s(v)=\E\tanh(Y_1(v))\tanh(Y_2(v)),
 \quad
 \begin{cases}
 \operatorname{Var}Y_1(v)=\operatorname{Var}Y_2(v)=\beta^2q,\\
 \operatorname{Cov}(Y_1(v),Y_2(v))=d+b^2v,
 \end{cases}
 \label{eq:signedf}
\end{equation}
and both Gaussian variables have mean $h$.  Gaussian differentiation
and Cauchy--Schwarz give
\begin{equation}
 f_s'(v)=b^2\E[\sech^2Y_1(v)\sech^2Y_2(v)]
 \le b^2a=s\alpha.
 \label{eq:signedcontraction}
\end{equation}
Since $f_s(q)=q$, it follows that
\begin{equation}
 f_s(v)-v\ge(1-s\alpha)(q-v)
 \qquad(-q\le v<q).
 \label{eq:signedslope}
\end{equation}
Thus the same Taylor argument gives a uniform gap on every closed
subset of $[-q,q)$.

For $v<-q$, the strict signed gap can be computed directly.  Write
$r=|v|$ and $t=b^2(r-q)$.  When $s>0$, one has $t>0$.  On $[r,1]$
the zero-source PDE factorizes, and since
\[
 \Psi(u,x)=\log\cosh x+\frac{b^2}{2}(1-u),
 \qquad q\le u\le1,
\]
its value at time $r$ is
\[
 U_s^{0,v}(r,x_1,x_2)
 =\log\cosh x_1+\log\cosh x_2+b^2(1-r).
\]
On $[q,r]$, $A_s^v=b^2\begin{psmallmatrix}1&-1\\-1&1\end{psmallmatrix}$
and $\gamma_v=1/2$.  Thus $W=\exp(U_s^{0,v}/2)$ solves the linear
backward heat equation in the direction $(1,-1)$, and
\begin{align}
 U_s^{0,v}(q,x_1,x_2)
 ={}&b^2(1-r)\nonumber\\
 &+2\log\E\sqrt{
   \cosh(x_1+\sqrt t\,Z)\cosh(x_2-\sqrt t\,Z)}.
 \label{eq:signedhopfcole}
\end{align}
The multiplier derivative passes through the same three levels
explicitly.  The $m=1$ level gives
\[
 W_s^v(r,x_1,x_2)=\tanh x_1\tanh x_2.
\]
Differentiating the $m=1/2$ semigroup in
\eqref{eq:signedhopfcole} gives
\[
 W_s^v(q,x_1,x_2)
 =
 \frac{\E\!\left[
 \sqrt{\cosh(x_1+\sqrt tZ)\cosh(x_2-\sqrt tZ)}
 \tanh(x_1+\sqrt tZ)\tanh(x_2-\sqrt tZ)\right]}
 {\E\sqrt{\cosh(x_1+\sqrt tZ)\cosh(x_2-\sqrt tZ)}}.
\]
Finally, the $m=0$ level below $q$ is ordinary Gaussian expectation:
\[
 \E W_s^v(0,Y_0,Y_0)
 =\E W_s^v(q,Y_0+b\sqrt q\,Z,Y_0-b\sqrt q\,Z).
\]
The Gaussian variables used at the successive levels are understood
to be independent.
Thus every zero-source multiplier identity used above is obtained by
successive differentiation of the $m=1$, $m=1/2$, and $m=0$
semigroups; no differentiability of an abstract weak PDE solution is
being assumed.

Cauchy--Schwarz and
$\E\cosh(x+\sqrt t\,Z)=e^{t/2}\cosh x$ give
\begin{equation}
 U_s^{0,v}(q,x_1,x_2)
 \le \Psi(q,x_1)+\Psi(q,x_2).
 \label{eq:signedCS}
\end{equation}
Equality in Cauchy--Schwarz would require
$\cosh(x_1+y)$ and $\cosh(x_2-y)$ to be proportional for every
$y\in\mathbb R$, which is equivalent to $x_1=-x_2$.

Below $q$, the PDE is linear.  Starting from $(Y_0,Y_0)$, its two
coordinates at time $q$ are
\[
 (Y_0+b\sqrt q\,Z,\;Y_0-b\sqrt q\,Z).
\]
Their sum is $2Y_0$.  Since $h>0$, one has
$\mathbb P(Y_0=0)=0$ (also when $d=0$, because then $Y_0=h$).
Therefore \eqref{eq:signedCS} is strict after averaging.  Its two
marginals both have the scalar local-field law, so
\begin{equation}
 \E U_s^{0,v}(0,Y_0,Y_0)<2\E\Psi(0,Y_0),
 \qquad v<-q, s>0.
 \label{eq:signedstrictPDE}
\end{equation}
The correction identity \eqref{eq:GTcorrection} now yields
\begin{equation}
 \mathfrak P_s(0,v)<2P_s^*,
 \qquad v<-q, s>0.
 \label{eq:signedpositivity}
\end{equation}
This is the signed GT positivity step in the present RS
normalization, derived without a Parisi free-energy formula.

We now prove the required uniformity, paying particular attention to
$s\downarrow0$.  At $s=0$ one has $b=0$ and
\[
 \mathfrak P_0(\lambda,v)
 =2\log2+\E f_\lambda(Y_0,Y_0)-\lambda v.
\]
The first derivative at zero is $q-v$, and the second derivative is
the variance, under the one-site coupled measure, of
$\varepsilon_1\varepsilon_2$; it is at most one.  Taylor's theorem
therefore gives, for every $\lambda\in\mathbb R$,
\begin{equation}
 \mathfrak P_0(\lambda,v)
 \le2P_0^*+\lambda(q-v)+\frac{\lambda^2}{2}.
 \label{eq:GTendpoint}
\end{equation}
For $v\le-q$, take $\lambda=v-q$.  By \eqref{eq:qcompact},
\begin{equation}
 \mathfrak P_0(v-q,v)
 \le2P_0^*-\frac12(v-q)^2
 \le2P_0^*-2q_{\cK}^2.
 \label{eq:GTsmallsgapzero}
\end{equation}
By Lemma~\ref{lem:GTcontinuity},
$\mathfrak P_s(\lambda,v)-2P_s^*$ is jointly continuous in
$(\beta,h,s,\lambda,v)$ on compact sets, including at $s=0$, $v=0$,
and $v=\pm q$.  Applying this continuity to
$\lambda=v-q$ and the compact set
\[
 \{(\beta,h,v):(\beta,h)\in\cK,\ -1\le v\le-q(\beta,h)\}
\]
shows that there exists $s_0>0$ such that
\begin{equation}
 \mathfrak P_s(v-q,v)\le2P_s^*-q_{\cK}^2,
 \qquad 0\le s\le s_0,\quad v\le-q.
 \label{eq:GTsmallsgap}
\end{equation}

It remains to treat $s\in[s_0,1]$.  At $v=-q$,
\eqref{eq:signedslope} gives
\[
 \partial_\lambda\mathfrak P_s(0,-q)
 \ge2\delta_{\cK}q_{\cK}.
\]
Joint continuity of this derivative from Lemma~\ref{lem:GTcontinuity} and
\eqref{eq:GTsecondderivative} yield $\eta_{\cK}>0$ such that, for
$-q-\eta_{\cK}\le v\le-q$ and $s\in[s_0,1]$,
\[
 \partial_\lambda\mathfrak P_s(0,v)
 \ge\delta_{\cK}q_{\cK}.
\]
The Cauchy--Schwarz argument above gives
$\mathfrak P_s(0,v)\le2P_s^*$ throughout this strip.  Taylor's theorem,
with
\[
 \ell_{\cK}=\min\left\{\lambda_{\cK},
 \frac{\delta_{\cK}q_{\cK}}{M_{\cK}}\right\},
 \qquad\lambda=-\ell_{\cK},
\]
therefore gives
\[
 \mathfrak P_s(-\ell_{\cK},v)
 \le2P_s^*-\frac12\delta_{\cK}q_{\cK}\ell_{\cK}.
\]
On the remaining compact set
\[
 \{(\beta,h,s,v):(\beta,h)\in\cK,\ s_0\le s\le1,\
                 -1\le v\le-q-\eta_{\cK}\},
\]
\eqref{eq:signedpositivity} is strict and continuous, so its minimum
gap is also positive.  Combining this with
\eqref{eq:GTsmallsgap} proves that, for some
$\kappa_{\cK}>0$,
\begin{equation}
 \inf_{\lambda\in\mathbb R}\mathfrak P_s(\lambda,v)
 \le2P_s^*-\kappa_{\cK},\qquad v\le-q,
 \label{eq:GTuniformnegativegap}
\end{equation}
uniformly on $\cK\times[0,1]$.

Finally, on $[0,1]$ away from $q$,
\eqref{eq:ATsign}, \eqref{eq:GTfirstderivative},
\eqref{eq:GTsecondderivative}, and compactness give a uniform positive
gap by choosing a sufficiently small multiplier of the opposite
sign.  Together with \eqref{eq:localGTgap},
\eqref{eq:signedslope}, and the uniform signed gap, this gives a
fixed positive gap away from $q$ and a quadratic gap near $q$.
Since $|v-q|\le2$, decreasing the constant if necessary yields
\eqref{eq:GTcoercivity} through the exact bound
\eqref{eq:GTfunctionalbound}.
\end{proof}

\begin{remark}
The additional parameter \eqref{eq:halfparameter} is the missing
parameter in the scalar Lata\l a interpolation.  If it is omitted, the
mixed derivative contains anomalous and longitudinal fluctuations.
With it, \eqref{eq:GTmixedderivative} is exactly
$-(1-s\alpha)$.
\end{remark}

\subsection{Finite-volume Gronwall closure}\label{subsec:gronwall}

Choose once and for all
\begin{equation}
 0<\lambda_0<c_{\cK}.
 \label{eq:lambda0}
\end{equation}
Define the quadratically coupled pressure and its normalized excess by
\begin{align}
 p^{(2)}_{N,s}(\lambda_0)
 &=
 \frac1{2N}\E\log
 \sum_{\sigma^1,\sigma^2}
 \exp\left\{
 H_{N,s}(\sigma^1)+H_{N,s}(\sigma^2)
 +\frac{\lambda_0N}{2}Q_{12}^2\right\},
 \label{eq:quadraticpressure}\\
 F_{N,s}(\lambda_0)
 &=
 p^{(2)}_{N,s}(\lambda_0)-\phi_N(s)
 =
 \frac1{2N}\E\log
 \left\langle e^{\lambda_0NQ_{12}^2/2}\right\rangle_s.
 \label{eq:normalizedcoupling}
\end{align}

\begin{lemma}[Sublinear coupled-pressure bound]
\label{lem:sublinearpressure}
With
\[
 \epsilon_N=C_{\cK}
 \sqrt{\frac{\log(N+1)}N}+\frac{C_{\cK}\log(N+1)}N,
\]
one has
\begin{equation}
 p^{(2)}_{N,s}(\lambda_0)\le P_s^*+\epsilon_N
 \label{eq:pressureenvelope}
\end{equation}
uniformly on $\cK\times[0,1]$.
\end{lemma}

\begin{proof}
Let $L_v=\log Z^{(2)}_{N,s}(v)$.  Each $L_v$ is a Lipschitz function
of the Gaussian disorder with squared Lipschitz constant at most
$C_{\cK}N$.  Indeed,
\[
 \left|\partial_{g_{ij}}L_v\right|
 \le\frac{2\beta\sqrt s}{\sqrt N},\qquad
 \left|\partial_{z_i}L_v\right|
 \le2\beta\sqrt{(1-s)q},
\]
and summing the squares over $i<j$ and $i$ gives the claim.  Hence
\begin{equation}
 \E\max_{v\in\mathcal R_N}(L_v-\E L_v)
 \le C_{\cK}\sqrt{N\log(N+1)}.
 \label{eq:gaussianmax}
\end{equation}
Using $|\mathcal R_N|=N+1$, \eqref{eq:GTcoercivity}, and
$\lambda_0<c_{\cK}$,
\begin{align*}
 2Np^{(2)}_{N,s}(\lambda_0)
 &=
 \E\log\sum_{v\in\mathcal R_N}
 \exp\left\{L_v+\frac{\lambda_0N}{2}(v-q)^2\right\}\\
 &\le
 \log(N+1)+
 \max_v\left\{\E L_v+\frac{\lambda_0N}{2}(v-q)^2\right\}
 +C_{\cK}\sqrt{N\log(N+1)}\\
 &\le2NP_s^*+2N\epsilon_N.
\end{align*}
This proves \eqref{eq:pressureenvelope}.
\end{proof}

Put
\begin{equation}
 D_N(s)=P_s^*-\phi_N(s).
 \label{eq:DN}
\end{equation}
At $s=0$, $D_N(0)=0$.  For $0<s<1$, Gaussian integration by parts
gives
\begin{equation}
 D_N'(s)=\frac{\beta^2}{4}A_s.
 \label{eq:DNprime}
\end{equation}
Both sides extend continuously to $s=0,1$ by finite-dimensional
Gaussian dominated convergence, so the identity also holds for the
one-sided endpoint derivatives.
On the other hand, convexity in the quadratic-coupling parameter gives
\[
 \frac14A_s
 =\left.\partial_\lambda F_{N,s}(\lambda)\right|_{\lambda=0+}
 \le\frac{F_{N,s}(\lambda_0)}{\lambda_0}.
\]
By \eqref{eq:pressureenvelope},
\[
 F_{N,s}(\lambda_0)
 \le D_N(s)+\epsilon_N.
\]
Consequently,
\begin{equation}
 D_N'(s)
 \le\frac{\beta^2}{\lambda_0}
 \bigl(D_N(s)+\epsilon_N\bigr).
 \label{eq:gronwallDN}
\end{equation}
Gronwall's inequality and compactness of $\cK$ imply
\begin{equation}
 \sup_{\cK\times[0,1]}D_N(s)\le C_{\cK}\epsilon_N,
 \qquad
 \sup_{\cK\times[0,1]}A_s\le C_{\cK}\epsilon_N.
 \label{eq:preconcentration}
\end{equation}
This is the first, deliberately nonoptimal, concentration estimate.
Its role is to normalize the coupled pressure without invoking the
Parisi formula.

\subsection{Fixed-deviation overlap concentration}\label{subsec:tails}

\begin{corollary}[Uniform fixed-deviation estimate]
\label{cor:tail}
For every $\varepsilon>0$, there are constants
$c_{\cK,\varepsilon},C_{\cK,\varepsilon}>0$ such that
\begin{equation}
 \sup_{\substack{(\beta,h)\in\cK\\0\le s\le1}}
 \nu_s\bigl(|Q_{12}|\ge\varepsilon\bigr)
 \le C_{\cK,\varepsilon}e^{-c_{\cK,\varepsilon}N}.
 \label{eq:tail}
\end{equation}
\end{corollary}

\begin{proof}
Let
\[
 Y_{N,s}
 =\log\left\langle
 e^{\lambda_0NQ_{12}^2/2}\right\rangle_s.
\]
By \eqref{eq:normalizedcoupling},
\eqref{eq:pressureenvelope}, and
\eqref{eq:preconcentration},
\[
 \E Y_{N,s}=2NF_{N,s}(\lambda_0)
 \le C_{\cK}N\epsilon_N=o(N)
\]
uniformly.  As a difference of one coupled and two uncoupled log
partition functions, $Y_{N,s}$ has squared Gaussian Lipschitz constant
at most $C_{\cK}N$.  Gaussian concentration therefore gives, for all
large enough $N$,
\[
 \mathbb P\left(
 Y_{N,s}>\frac{\lambda_0\varepsilon^2N}{4}\right)
 \le C_{\cK,\varepsilon}e^{-c_{\cK,\varepsilon}N}.
\]
Outside this exceptional event,
\[
 \left\langle\1_{\{|Q_{12}|\ge\varepsilon\}}\right\rangle_s
 \le
 \exp\left\{-\frac{\lambda_0\varepsilon^2N}{2}
 +Y_{N,s}\right\}
 \le e^{-\lambda_0\varepsilon^2N/4}.
\]
Taking disorder expectation and enlarging the constant for the
finitely many remaining $N$ proves \eqref{eq:tail}.
\end{proof}

\subsection{The finite-size cavity system}\label{subsec:cavity}

The next proposition is the one-species cavity fluctuation system.  We
include its derivation because the size and form of its remainder are
essential for the absorption step.

\begin{proposition}[Three-dimensional cavity equations]
\label{prop:cavity}
Let $x_s=(A_s,B_s,C_s)^{\mathsf T}$.  Uniformly on
$\cK\times[0,1]$,
\begin{equation}
 (I-sL)x_s=\frac1N\theta+\rho_{N,s},
 \label{eq:cavitysystem}
\end{equation}
where
\begin{equation}
 \theta=
 \begin{pmatrix}
 1-q^2\\ q-q^2\\ r-q^2
 \end{pmatrix},
 \label{eq:theta}
\end{equation}
\begin{equation}
 L=
 \begin{pmatrix}
 b_2&-4b_1&3b_0\\
 b_1&b_2-2b_1-3b_0&6b_0-3b_1\\
 b_0&4b_1-8b_0&b_2-8b_1+10b_0
 \end{pmatrix},
 \label{eq:L}
\end{equation}
with
\begin{equation}
 b_2=\beta^2(1-q^2),\qquad
 b_1=\beta^2q(1-q),\qquad
 b_0=\beta^2(r-q^2),
 \label{eq:bs}
\end{equation}
and
\begin{equation}
 \norm{\rho_{N,s}}_\infty
 \le C_{\cK}
 \left(N^{-3/2}+\nu_s|Q_{12}|^3\right).
 \label{eq:rho}
\end{equation}
\end{proposition}

\begin{proof}
Write $\varepsilon_a=\sigma_N^a$ and
\[
 Q^-_{ab}=\frac1N\sum_{i<N}\sigma_i^a\sigma_i^b-q.
\]
By site exchangeability after disorder averaging,
\begin{align}
 A_s
 &=\nu_s\bigl[Q_{12}(\varepsilon_1\varepsilon_2-q)\bigr],\nonumber\\
 B_s
 &=\nu_s\bigl[Q_{12}(\varepsilon_1\varepsilon_3-q)\bigr],\nonumber\\
 C_s
 &=\nu_s\bigl[Q_{12}(\varepsilon_3\varepsilon_4-q)\bigr].
 \label{eq:siteexchange}
\end{align}
Since
$Q_{12}=Q^-_{12}+N^{-1}\varepsilon_1\varepsilon_2$,
each line consists of a cavity term involving $Q^-_{12}$ and a
diagonal term of order $N^{-1}$; explicitly,
\begin{align}
 A_s&=\nu_s[Q^-_{12}(\varepsilon_1\varepsilon_2-q)]
 +\frac1N\nu_s[\varepsilon_1\varepsilon_2
                    (\varepsilon_1\varepsilon_2-q)],\nonumber\\
 B_s&=\nu_s[Q^-_{12}(\varepsilon_1\varepsilon_3-q)]
 +\frac1N\nu_s[\varepsilon_1\varepsilon_2
                    (\varepsilon_1\varepsilon_3-q)],\nonumber\\
 C_s&=\nu_s[Q^-_{12}(\varepsilon_3\varepsilon_4-q)]
 +\frac1N\nu_s[\varepsilon_1\varepsilon_2
                    (\varepsilon_3\varepsilon_4-q)].
 \label{eq:cavitydecomposition}
\end{align}

We define the last-spin interpolation explicitly.  Write
$\sigma=(\rho,\varepsilon)$ with
$\rho\in\{-1,1\}^{N-1}$ and let
\begin{align*}
 H^-_{N,s}(\rho)
 ={}&\frac{\beta\sqrt s}{\sqrt N}
 \sum_{1\le i<j<N}g_{ij}\rho_i\rho_j
 +\sum_{i<N}\bigl(h+\beta\sqrt{(1-s)q}\,z_i\bigr)\rho_i.
\end{align*}
Let $\widehat z$ be a standard Gaussian independent of all existing
disorder.  For $0\le u\le1$, set
\begin{align}
 H_{N,s,u}(\rho,\varepsilon)
 ={}&H^-_{N,s}(\rho)
 +\bigl(h+\beta\sqrt{(1-s)q}\,z_N\bigr)\varepsilon\nonumber\\
 &+\frac{\beta\sqrt{su}}{\sqrt N}
   \sum_{i<N}g_{iN}\rho_i\varepsilon
 +\beta\sqrt{s(1-u)q}\,\widehat z\,\varepsilon .
 \label{eq:cavityHamiltonian}
\end{align}
Denote its disorder--Gibbs expectation by $\nu_{s,u}$.  At $u=1$,
\eqref{eq:cavityHamiltonian} is exactly \eqref{eq:path}, so
$\nu_{s,1}=\nu_s$.  At $u=0$, the last spin is independent of the
first $N-1$ spins and sees the Gaussian field
\[
 Y\stackrel d=h+\beta\sqrt q\,Z.
\]
Conditional on $Y$, its replicas are independent, with mean
$\tanh Y$.  Thus, at the decoupled endpoint,
\begin{equation}
 \E_0(\varepsilon_a\varepsilon_b)=q,\qquad
 \E_0(\varepsilon_a\varepsilon_b\varepsilon_c\varepsilon_d)=r
 \label{eq:epsilonmoments}
\end{equation}
for distinct indices.

For completeness, we differentiate this interpolation.  The
derivative of the covariance between the last-spin parts of two
configurations $(\rho,\varepsilon)$ and
$(\rho',\varepsilon')$ is
\begin{equation}
 s\beta^2\varepsilon\varepsilon'
 \left(\frac1N\sum_{i<N}\rho_i\rho_i'-q\right)
 =s\beta^2\varepsilon\varepsilon'Q^-(\rho,\rho').
 \label{eq:cavitycovariancederivative}
\end{equation}
For a bounded function $F$ of $n$ replicas, differentiating the
normalized Gibbs expectation for $0<u<1$ and applying Gaussian
integration by parts to \eqref{eq:cavitycovariancederivative} gives
\begin{align}
 \frac{\dd}{\dd u}\nu_{s,u}(F)
 =s\beta^2\Bigg\{&
 \sum_{1\le a<b\le n}
 \nu_{s,u}(F\varepsilon_a\varepsilon_bQ^-_{ab})
 -n\sum_{a=1}^n
 \nu_{s,u}(F\varepsilon_a\varepsilon_{n+1}Q^-_{a,n+1})
 \nonumber\\
 &+\frac{n(n+1)}2
 \nu_{s,u}(F\varepsilon_{n+1}\varepsilon_{n+2}
 Q^-_{n+1,n+2})\Bigg\}.
 \label{eq:cavityderivative}
\end{align}
Both sides of \eqref{eq:cavityderivative} extend continuously to
$u=0,1$.  We use those extensions for the one-sided endpoint
derivatives below; thus the square-root coefficients in
\eqref{eq:cavityHamiltonian} cause no endpoint singularity.
To see all coefficients in one formula, define the replica operator
\begin{align*}
 \mathcal D_nF={}&
 \sum_{1\le a<b\le n}
 F\varepsilon_a\varepsilon_bQ^-_{ab}
 -n\sum_{a=1}^n
 F\varepsilon_a\varepsilon_{n+1}Q^-_{a,n+1}\\
 &+\frac{n(n+1)}2
 F\varepsilon_{n+1}\varepsilon_{n+2}Q^-_{n+1,n+2}.
\end{align*}
Then \eqref{eq:cavityderivative} is simply
$\frac{\dd}{\dd u}\nu_{s,u}(F)
=s\beta^2\nu_{s,u}(\mathcal D_nF)$.  This form also makes the
second differentiation below unambiguous: each application of
$\mathcal D$ appends exactly one centered cavity-overlap factor, and
the replica number is increased only to include the displayed new
indices.

We first record the comparison estimate that permits us to replace all
intermediate cavity expectations by the full endpoint expectation.  By
construction, $\nu_{s,1}=\nu_s$.  Set
\begin{equation}
 M_3(u)=\nu_{s,u}|Q^-_{12}|^3.
 \label{eq:cavityM3}
\end{equation}
Applying \eqref{eq:cavityderivative} with
$F=|Q^-_{12}|^3$, using $|Q^-_{ab}|\le2$ and replica symmetry, gives
\begin{equation}
 |M_3'(u)|\le C_{\cK}M_3(u).
 \label{eq:cavityM3derivative}
\end{equation}
Indeed, every term on the right is bounded by
$2\nu_{s,u}|Q^-_{12}|^3$.  Integrating backward from $u=1$ and applying
Gronwall's inequality yields
\begin{equation}
 \sup_{0\le u\le1}M_3(u)\le C_{\cK}M_3(1).
 \label{eq:cavityM3comparison}
\end{equation}
Since
$Q^-_{12}=Q_{12}-N^{-1}\varepsilon_1\varepsilon_2$ and
$|\varepsilon_1\varepsilon_2|=1$, the inequality
$(x+y)^3\le4(x^3+y^3)$ for $x,y\ge0$ gives
\begin{equation}
 M_3(1)\le C\left(\nu_s|Q_{12}|^3+N^{-3}\right).
 \label{eq:cavityM3endpoint}
\end{equation}

Next let $P^-$ be any product of two centered cavity overlaps, for
example $(Q^-_{12})^2$, $Q^-_{12}Q^-_{13}$, or
$Q^-_{12}Q^-_{34}$.  Differentiating $\nu_{s,u}(P^-)$ by
\eqref{eq:cavityderivative} produces products of three centered cavity
overlaps, possibly multiplied by spin variables.  By H\"older's
inequality and replica symmetry, each such term is bounded by $M_3(u)$.
Consequently,
\begin{equation}
 \left|\nu_{s,u}(P^-)-\nu_{s,v}(P^-)\right|
 \le C_{\cK}|u-v|M_3(1),
 \qquad 0\le u,v\le1.
 \label{eq:cavityquadraticcomparison}
\end{equation}
At the full endpoint, replacing either factor $Q^-_{ab}$ by $Q_{ab}$
costs at most
$C(N^{-1}|Q_{cd}|+N^{-2})$.  The pointwise Young inequality
\begin{equation}
 \frac{|x|}{N}\le\frac{|x|^3}{3}+\frac{2}{3N^{3/2}}
 \label{eq:cavityYoung}
\end{equation}
therefore implies, for each of the three quadratic overlap patterns,
\begin{equation}
 \left|\nu_{s,0}(P^-)-\nu_s(P)\right|
 \le C_{\cK}\left(\nu_s|Q_{12}|^3+N^{-3/2}\right),
 \label{eq:cavityendpointcomparison}
\end{equation}
where $P$ denotes the corresponding product of full centered overlaps.

Apply the first-order Taylor formula at $u=0$ to the three cavity
terms in \eqref{eq:cavitydecomposition}.  Their zeroth-order terms vanish by
\eqref{eq:epsilonmoments}; thus, for each such function $F$,
\[
 \nu_{s,1}(F)=
 \left.\frac{\dd}{\dd u}\nu_{s,u}(F)\right|_{u=0}
 +\int_0^1(1-u)\frac{\dd^2}{\dd u^2}\nu_{s,u}(F)\,\dd u.
\]
Formula \eqref{eq:cavityderivative} first
expresses the derivatives at $u=0$ in terms of the decoupled quadratic
moments.  Replacing those moments by $A_s,B_s,C_s$ using
\eqref{eq:cavityendpointcomparison} changes each equation by at most
$C_{\cK}(\nu_s|Q_{12}|^3+N^{-3/2})$.  Modulo this controlled error, the
first-order coefficients are as follows:
\begin{equation}
\begin{array}{c|ccc|c}
 &A_s&B_s&C_s&N^{-1}\text{ term}\\ \midrule
 A_s&b_2&-4b_1&3b_0&N^{-1}(1-q^2)\\
 B_s&b_1&b_2-2b_1-3b_0&6b_0-3b_1&N^{-1}(q-q^2)\\
 C_s&b_0&4b_1-8b_0&b_2-8b_1+10b_0&N^{-1}(r-q^2).
\end{array}
\label{eq:cavitytable}
\end{equation}
The factor $s$ from \eqref{eq:cavityderivative} is not included in
the table; it is the common factor producing $I-sL$ in
\eqref{eq:cavitysystem}.
For example, the first line uses $n=2$.  Before replacing decoupled
quadratic moments by endpoint moments, the three terms in
\eqref{eq:cavityderivative} are
\[
 (1-q^2)\nu_{s,0}((Q^-_{12})^2),\qquad
 -4q(1-q)\nu_{s,0}(Q^-_{12}Q^-_{13}),\qquad
 3(r-q^2)\nu_{s,0}(Q^-_{12}Q^-_{34}).
\]
Equation \eqref{eq:cavityendpointcomparison} turns these into the first
row of \eqref{eq:cavitytable}, with an error of the required size.
For the second line take $n=3$ and
$F=(\varepsilon_1\varepsilon_3-q)Q^-_{12}$.  Grouping the pairs in the
first sum of \eqref{eq:cavityderivative} according to whether they are
$12$, $13$, or $23$, and doing the same for the terms involving the
new replica, gives
\[
 b_1A_s+(b_2-2b_1-3b_0)B_s+(6b_0-3b_1)C_s.
\]
For the third line take $n=4$ and
$F=(\varepsilon_3\varepsilon_4-q)Q^-_{12}$.  The six old-replica
pairs, the four terms containing replica $5$, and the final
$5,6$ term give, respectively,
\[
 b_0A_s+(4b_1-8b_0)B_s+(b_2-8b_1+10b_0)C_s.
\]
Here is a direct rule that verifies every entry.  Conditional on $Y$,
the $\varepsilon_a$ are independent with mean $\tanh Y$.  Hence a
spin monomial has expectation $1$, $q$, or $r$ according as the
number of replica indices occurring an odd number of times is
$0$, $2$, or $4$.  In particular, if $e$ is a replica edge, then
\[
 \E_0[(\varepsilon_p\varepsilon_q-q)\varepsilon_e]
 =
 \begin{cases}
  1-q^2,&e=\{p,q\},\\
  q-q^2,&e\text{ shares exactly one endpoint with }\{p,q\},\\
  r-q^2,&e\text{ is disjoint from }\{p,q\}.
 \end{cases}
\]
Counting the old, one-new, and two-new edges in $\mathcal D_2$,
$\mathcal D_3$, and $\mathcal D_4$ gives respectively
\[
 (b_2,-4b_1,3b_0),
\]
\[
 (b_1,\ b_2-2b_1-3b_0,\ 6b_0-3b_1),
\]
and
\[
 (b_0,\ 4b_1-8b_0,\ b_2-8b_1+10b_0),
\]
which verifies all three rows of \eqref{eq:cavitytable}.

The diagonal terms come from
$N^{-1}\varepsilon_1\varepsilon_2$ in the exact identity
$Q_{12}=Q^-_{12}+N^{-1}\varepsilon_1\varepsilon_2$.  For example,
\[
 \E_0[\varepsilon_1\varepsilon_2
       (\varepsilon_1\varepsilon_3-q)]
 =\E\tanh^2Y-q^2=q-q^2.
\]
The other two diagonal expectations are
$\E_0[\varepsilon_1\varepsilon_2
(\varepsilon_1\varepsilon_2-q)]=1-q^2$ and
$\E_0[\varepsilon_1\varepsilon_2
(\varepsilon_3\varepsilon_4-q)]=r-q^2$.

It remains to bound the integral Taylor remainder.  Each of the three
cavity functions is a bounded spin factor times one centered cavity
overlap.  The first application of $\mathcal D_n$ gives a finite sum
of bounded spin factors times two centered overlaps.  Applying
\eqref{eq:cavityderivative} once more, now with all replica indices
appearing in each summand included, gives a finite sum of terms
\[
 c\,\nu_{s,u}\bigl[
 \chi(\varepsilon)\,Q^-_{a_1b_1}Q^-_{a_2b_2}Q^-_{a_3b_3}\bigr],
\]
where $|c|\le C$ and $|\chi|\le C$; the number of terms is bounded
because the original replica numbers are only $2$, $3$, and $4$.
H\"older's inequality,
replica symmetry, \eqref{eq:cavityM3comparison}, and
\eqref{eq:cavityM3endpoint} give, uniformly in $u$,
\begin{equation}
 \left|\frac{\dd^2}{\dd u^2}\nu_{s,u}(F)\right|
 \le C_{\cK}\left(\nu_s|Q_{12}|^3+N^{-3}\right)
 \label{eq:cavitysecondderivativebound}
\end{equation}
for each of the three cavity functions $F$ occurring in
\eqref{eq:cavitydecomposition}.  Hence their integral Taylor remainders obey
the same bound.

For the $N^{-1}$ diagonal terms, one differentiation gives a bounded
spin observable times one centered cavity overlap.  Thus their change
between $u=0$ and $u=1$ is bounded by
\[
 \frac{C_{\cK}}N\sup_{0\le u\le1}\nu_{s,u}|Q^-_{12}|.
\]
Using \eqref{eq:cavityYoung} pointwise and then
\eqref{eq:cavityM3comparison}--\eqref{eq:cavityM3endpoint}, this is at
most
\begin{equation}
 C_{\cK}\left(\nu_s|Q_{12}|^3+N^{-3/2}\right).
 \label{eq:cavitydiagonalerror}
\end{equation}
Combining \eqref{eq:cavityendpointcomparison},
\eqref{eq:cavitysecondderivativebound}, and
\eqref{eq:cavitydiagonalerror} proves \eqref{eq:rho}.  The coefficient
table \eqref{eq:cavitytable} is therefore equivalent to
\eqref{eq:cavitysystem} with the stated remainder.
\end{proof}

\subsection{Stability of the cavity matrix}\label{subsec:stability}

Put
\begin{equation}
 \gamma'=1-4q+3r,\qquad
 \gamma''=2q+q^2-3r.
 \label{eq:gammas}
\end{equation}
The change of coordinates
\[
 S=
 \begin{pmatrix}
 0&2&-3\\
 1&-4&3\\
 1&-2&1
 \end{pmatrix}
\]
gives, by direct multiplication,
\begin{equation}
 SLS^{-1}
 =
 \begin{pmatrix}
 \beta^2\gamma'&\beta^2\gamma''&0\\
 0&\beta^2\gamma'&0\\
 0&0&\alpha
 \end{pmatrix}.
 \label{eq:triangular}
\end{equation}
Hence
\begin{equation}
 \det(I-sL)
 =(1-s\alpha)(1-s\beta^2\gamma')^2.
 \label{eq:det}
\end{equation}
Since $r\le q$,
\[
 \alpha-\beta^2\gamma'
 =2\beta^2(q-r)\ge0.
\]
Thus $\beta^2\gamma'\le\alpha<1$.  If $\gamma'<0$, then
$1-s\beta^2\gamma'\ge1$; if $\gamma'\ge0$, then
$1-s\beta^2\gamma'\ge1-\alpha$.  Together with
\eqref{eq:delta}, this gives
\begin{equation}
 M_{\cK}:=
 \sup_{\substack{(\beta,h)\in\cK\\0\le s\le1}}
 \norm{(I-sL)^{-1}}_\infty<\infty.
 \label{eq:inverse}
\end{equation}
Notice that \eqref{eq:triangular}, rather than the eigenvalues alone,
also controls the possible Jordan term.

\subsection{Absorption of the cubic remainder}\label{subsec:absorb}

Replica symmetry and Cauchy--Schwarz give
\[
 |B_s|\le A_s,\qquad |C_s|\le A_s.
\]
Therefore $\norm{x_s}_\infty=A_s$.  Inverting
\eqref{eq:cavitysystem}, and using \eqref{eq:inverse} and
\eqref{eq:rho}, yields
\begin{equation}
 A_s
 \le \frac{C_{\cK}}N
 +C_{\cK}N^{-3/2}
 +C_{\cK}\nu_s|Q_{12}|^3.
 \label{eq:preabsorb}
\end{equation}

Fix $\varepsilon>0$.  Since $|Q_{12}|\le2$,
\begin{align}
 \nu_s|Q_{12}|^3
 &\le
 \varepsilon\,\nu_s(Q_{12}^2)
 +8\,\nu_s(|Q_{12}|>\varepsilon)\nonumber\\
 &\le
 \varepsilon A_s
 +8C_{\cK,\varepsilon}e^{-c_{\cK,\varepsilon}N},
 \label{eq:thirdsplit}
\end{align}
where \cref{cor:tail} was used.  Choose $\varepsilon=\varepsilon_{\cK}$
so that the coefficient of $A_s$ obtained by substituting
\eqref{eq:thirdsplit} into \eqref{eq:preabsorb} is at most $1/2$.
We obtain
\[
 A_s\le\frac{C_{\cK}}N
\]
uniformly on $\cK\times[0,1]$, after increasing $C_{\cK}$ to include
the finitely many small values of $N$.  This proves the first assertion
of \cref{thm:main}.

\subsection{Free energy}\label{subsec:freeenergy}

Recall $\phi_N(s)$ from \eqref{eq:pathfreeenergy}.  For $0<s<1$,
differentiating under the expectation and applying \eqref{eq:GIBP} to
every $g_{ij}$ and $z_i$ gives
\begin{align*}
 \phi_N'(s)
 ={}&\frac{\beta^2}{2N^2}\sum_{i<j}
 \E\left[1-\langle\sigma_i\sigma_j\rangle_s^2\right]
 -\frac{\beta^2q}{2N}\sum_i
 \E\left[1-\langle\sigma_i\rangle_s^2\right]\\
 ={}&\frac{\beta^2}{4}
 \left(1-2q+q^2-\nu_s(R_{12}-q)^2\right).
\end{align*}
In the second line we used
\[
 \sum_{i<j}\sigma_i^1\sigma_j^1\sigma_i^2\sigma_j^2
 =\frac12\left(N^2R_{12}^2-N\right)
\]
and
$N^{-1}\sum_i\E\langle\sigma_i^1\sigma_i^2\rangle_s
=\nu_s(R_{12})$.  Thus
\begin{equation}
 \phi_N'(s)=\frac{\beta^2}{4}
 \left((1-q)^2-\nu_s(Q_{12}^2)\right).
 \label{eq:freeenergyderivative}
\end{equation}
The right-hand side is continuous at $s=0,1$ by finite-dimensional
Gaussian dominated convergence.  Hence
\eqref{eq:freeenergyderivative} extends to the one-sided endpoint
derivatives.

At $s=0$ the Gibbs measure is a product measure, and
\[
 \phi_N(0)=\log2+\E\log\cosh(h+\beta\sqrt q\,Z).
\]
Integrating \eqref{eq:freeenergyderivative} gives
\begin{equation}
 \Phi_{\RS}(\beta,h)-\phi_N(1)
 =
 \frac{\beta^2}{4}\int_0^1\nu_s(Q_{12}^2)\,\dd s.
 \label{eq:freeenergyidentity}
\end{equation}
The integrand is nonnegative, and the already proved second-moment
bound makes the right-hand side at most $C_{\cK}/N$.  This proves the
free-energy assertion of \cref{thm:main}.

\subsection{Replicon susceptibility}\label{subsec:susceptibility}

Define
\[
 D_s=A_s-2B_s+C_s.
\]
The last row of the transformed system \eqref{eq:triangular}, or
equivalently the scalar product of \eqref{eq:cavitysystem} with
$(1,-2,1)$, gives
\begin{equation}
 (1-s\alpha)D_s=\frac aN+\rho^{(R)}_{N,s},
 \label{eq:repliconeq}
\end{equation}
where
\begin{equation}
 |\rho^{(R)}_{N,s}|
 \le C_{\cK}\left(N^{-3/2}+\nu_s|Q_{12}|^3\right).
 \label{eq:repliconrho}
\end{equation}
Here
\[
 (1,-2,1)\theta=1-2q+r=a.
\]

The $O(N^{-1})$ second-moment estimate and
\eqref{eq:thirdsplit} imply, for every fixed $\varepsilon>0$,
\[
 \sup_{\cK\times[0,1]}
 N\nu_s|Q_{12}|^3
 \le C_{\cK}\varepsilon
 +8N C_{\cK,\varepsilon}e^{-c_{\cK,\varepsilon}N}.
\]
First let $N\to\infty$ and then $\varepsilon\downarrow0$.  Thus
\[
 \nu_s|Q_{12}|^3=o_{\cK}(N^{-1})
\]
uniformly along the path.  It follows from
\eqref{eq:repliconeq}, \eqref{eq:repliconrho}, and
$1-s\alpha\ge\delta_{\cK}$ that
\[
 ND_s
 =
 \frac{a}{1-s\alpha}+o_{\cK}(1),
\]
which is the final assertion of \cref{thm:main}.

\begin{thebibliography}{99}

\bibitem{ChenGT}
W.-K.~Chen,
\emph{Variational representations for the Parisi functional and the
two-dimensional Guerra--Talagrand bound},
Ann. Probab. \textbf{45} (2017), 3929--3966.
\href{https://arxiv.org/abs/1501.06635}{arXiv:1501.06635}.

\bibitem{DeyWu}
P.~S.~Dey and Q.~Wu,
\emph{Fluctuation results for multi-species
Sherrington--Kirkpatrick model in the replica symmetric regime},
J. Stat. Phys. \textbf{185} (2021), Paper No.~22.
\href{https://arxiv.org/abs/2012.13381}{arXiv:2012.13381}.

\bibitem{Lopatto}
P.~Lopatto,
\emph{Replica symmetry up to the de Almeida--Thouless line in the
Sherrington--Kirkpatrick model},
preprint (2026).
\href{https://arxiv.org/abs/2604.11921}{arXiv:2604.11921}.

\bibitem{Panchenko}
D.~Panchenko,
\emph{The Sherrington--Kirkpatrick Model},
Springer Monographs in Mathematics, Springer, 2013.

\bibitem{Talagrand}
M.~Talagrand,
\emph{Mean Field Models for Spin Glasses}, Vols.~I--II,
Springer, 2011.

\end{thebibliography}


\end{document}
